\section{Introduction}
\label{sec:introduction}

Deep neural networks are typically evaluated by predictive accuracy, yet deployment safety also depends on the structure of the representation a network learns internally, not only on how it classifies a fixed test set. Whether downstream reliability estimation succeeds therefore depends not only on what a representation encodes, but also on how that information is organized. A representation's geometry---how tightly points cluster by class, how well-conditioned the resulting covariance is, how closely that structure resembles the class-conditional Gaussian shape many downstream estimators assume---is itself shaped by training, and can change substantially without moving classification accuracy at all.

Distance-based reliability and out-of-distribution (OOD) estimation methods are built directly on this geometry. Mahalanobis distance \citep{lee2018simple} scores a test point by its distance to the nearest class-conditional mean under an estimated covariance; cosine-similarity-to-centroid variants replace the covariance-weighted distance with a simpler angular one; non-parametric $k$-nearest-neighbor scorers \citep{sun2022out} drop the notion of a class centroid altogether and use local density in the representation space instead. Despite spanning fully parametric to fully non-parametric, all three rest on the same premise: whatever geometric structure a representation has, a sufficiently faithful measure of distance within it will track how typical, and therefore how trustworthy, a given input is. Other reliability estimators relax this purely geometric premise---Energy \citep{liu2020energy} and Virtual-Logit Matching \citep{wang2022vim} incorporate the classifier's own logits directly, and density-based estimators replace a single Gaussian assumption with a non-parametric one---but share the same underlying constraint: all are computed from a fixed, already-trained representation, without retraining or access to the original objective.

This premise ties reliability estimation to representation geometry by construction, so any training procedure that reshapes representation geometry also reshapes, whether or not this is intended, the reliability estimator built on top of it. Domain-adversarial and orthogonality-based disentanglement training \citep{ganin2015unsupervised} are procedures designed explicitly to reorganize this geometry---separating or suppressing structure associated with a nuisance variable such as acquisition site or imaging domain---and are increasingly used in clinical imaging pipelines for exactly that reason. Whether the geometric changes such training induces are the kind distance-based reliability estimation can still read, or the kind that severs the link between geometry and reliability without affecting classification accuracy, has not, to our knowledge, been audited directly.

We audit this assumption using a disentanglement dose-response ladder that isolates the effect of orthogonality strength while holding the architecture and representation fixed. This design allows us to ask whether geometric change, distance-based reliability estimation, and information decodability remain coupled under disentanglement. Doing so lets us distinguish two claims the existing OOD literature does not, to our knowledge, separate for a training-induced geometric shift: that a representation has lost the information a reliability estimator would need, versus that the representation still holds this information in a form distance-based geometry cannot use.

Across the ladder, representation geometry changes substantially, whereas eight structurally distinct reliability estimators---five distance-based, three not---consistently remain near or below chance ($\sim$0.40 AUROC), with a single exception discussed as a hypothesis-generating observation rather than a counter-finding (Section~\ref{sec:discussion}). A supervised probe applied to the identical representations, in contrast, recovers the relevant domain signal at 0.73--0.81 AUROC. These results are obtained within a domain-adversarial training regime evaluated against its own adversarial target domain (Section~\ref{sec:datasets}), and are scoped accordingly throughout this paper. Information can remain decodable while becoming largely inaccessible to non-probing reliability estimators.
